\tableofcontents

\listoffigures
\listoftables
\lstlistoflistings

\chapter{Dokumentinformationen}

% resets page numbers to 1 and continues in arabic
\pagenumbering{arabic}
\setcounter{page}{1}

\begin{table}[h]
	\begin{tabularx}{\textwidth}{|l|X|}
		\hline
		\textbf{Datum} & 21.01.26\\
		\hline
		\textbf{Version} & 0.6\\
		\hline 
		\textbf{Verantwortlicher Autor} & Nils Bennerscheidt\\
		\hline 
	\end{tabularx}
	\label{tbl:dokumentinformation}
	\captionof{table}{Dokumentinformationen}
\end{table}

\section{Dokumenthistorie}

\begin{table}[h]
	\begin{tabularx}{\textwidth}{|X|X|X|X|}
		\hline
		\textbf{Datum} & \textbf{Version} & \textbf{Änderung} & \textbf{Bearbeiter}\\
		\hline
		08.08.25 & 0.1 & Initiale Erstellung & Frank Sieger\\
		\hline
		15.08.25 & 0.2 & Erweiterung Asynchronität und Attribute & Frank Sieger\\
		\hline
		17.10.25 & 0.3 & Erweiterung Logging Web-Client & Lukas Stucki\\
		\hline
		23.11.25 & 0.4.1 & Umstrukturierung und Erweiterung auf Dokumentation aller LO-Services  & Nils Bennerscheidt\\
		\hline
		16.01.26 & 0.5 & Update Template  & Bastian Schuchardt\\
		\hline
		20.01.26 & 0.5.1 & Font   & Nils Bennerscheidt\\
		\hline
		21.01.26 & 0.5.2 & Neue SNDE Logs   & Nils Bennerscheidt\\
		\hline
		21.01.26 & 0.5.3 & Template Korrektur & Bastian Schuchardt \\
		\hline
		22.01.26 & 0.5.4 & ToDos, Neuere Logs und Kleinigkeiten & Nils Bennerscheidt \\
		\hline
		17.03.26 & 0.5.5 & Überarbeitung & Nils Bennerscheidt \\
		\hline
	\end{tabularx}
	\label{tbl:dokumenthistorie}
	\captionof{table}{Dokumenthistorie}
		
\end{table}

\chapter{Präambel}Ziel dieses Dokuments ist es, eine Übersicht über alle Logfiles, die von mLogistics erstellt werden, zu bieten und darüber hinaus eine sinnvolle Ordnung in der Ordner- und Dateistruktur darzustellen. Es dient als Grundlage für den Betrieb, die Analyse sowie die Weiterentwicklung der Logging-Strategie innerhalb der Systemlandschaft.

\section{Rahmenbedingungen}
Dieses Dokument ist noch nicht final ausgearbeitet; Änderungen und Ergänzungen sind vorbehalten.
Es beschreibt die Konfigurationen und Umgebungen, wie sie von LogObject empfohlen werden und für einen stabilen Betrieb vorgesehen sind.

\section{Zielgruppe}

Dieses Dokument richtet sich primär an Applikationsverantwortliche der mLogistics-Systemlandschaft. Dazu zählen insbesondere Personen, die für den Betrieb, die Konfiguration sowie die Analyse der Anwendungen und deren Logging verantwortlich sind.

Darüber hinaus adressiert dieses Dokument auch indirekt alle beteiligten Rollen, die im erweiterten Betrieb und Support der Systeme tätig sind. Hierzu gehören beispielsweise Systemadministratoren, DevOps-Teams sowie externe Dienstleister, die mit Wartung, Fehleranalyse oder Infrastrukturbetreuung beauftragt sind.

Ziel ist es, sowohl eine fachliche als auch eine technische Grundlage bereitzustellen, um:
\begin{itemize}
\item die Struktur und Funktionsweise des Loggings zu verstehen,
\item relevante Logdateien schnell identifizieren zu können,
\item Fehler effizient zu analysieren und einzugrenzen,
\item sowie den stabilen Betrieb der Systeme sicherzustellen.
\end{itemize}

Das Dokument setzt ein grundlegendes technisches Verständnis von verteilten Systemen und Java-basierten Anwendungen voraus, ist jedoch so aufgebaut, dass auch angrenzende Rollen ohne tiefgehende Entwicklungskenntnisse die wesentlichen Inhalte nachvollziehen können.
